% © 2026 Ing. David Jaroš — CC BY-NC-ND 4.0
%
% File: research_tracks/EW/weinberg_angle_rg_twisted_scale_closure.tex
% Status: RESEARCH TRACK / L1 CONDITIONAL.
% Purpose: RG descent from the UBT Weinberg boundary sin^2(theta_W)=3/8 to M_Z.

\documentclass[12pt,a4paper]{article}
\usepackage[a4paper,margin=2.5cm]{geometry}
\usepackage{amsmath,amssymb,mathtools,booktabs}
\usepackage[most]{tcolorbox}
\usepackage{hyperref}
\usepackage{xcolor}
\usepackage{microtype}

\hypersetup{colorlinks=true,linkcolor=blue!60!black,urlcolor=blue!60!black,citecolor=green!50!black}

\newtcolorbox{statusbox}[1][]{
  colback=yellow!8!white,
  colframe=orange!70!black,
  arc=3pt,
  boxrule=1pt,
  title=Status,
  #1
}

\title{RG Descent of the UBT Weinberg Boundary}
\author{Ing. David Jaroš}
\date{2026-06-14}

\begin{document}
\maketitle

\begin{abstract}
The companion theorem \texttt{weinberg\_angle\_gut\_norm\_theorem.tex}
proves the UBT boundary condition
\[
  \sin^2\theta_W(M_{\rm UBT})=\frac38
\]
from hypercharge and weak-isospin generator norms.  This note attacks the
second step: RG descent to the $Z$ pole.  Minimal non-supersymmetric one-loop
SM running does not reproduce the observed value from the $3/8$ boundary by
itself.  However, the odd-winding spinor sector already used in the alpha proof
naturally points to a twisted Layer2 threshold with MSSM-like one-loop beta
coefficients
\[
  b_1=\frac{33}{5},\qquad b_2=1.
\]
With $M_{\rm UBT}\simeq2\times10^{16}\,{\rm GeV}$ and
$\alpha_{\rm UBT}^{-1}\simeq24.3$, this gives
\[
  \sin^2\theta_W(M_Z)=0.23114\ldots,
\]
within $8\times10^{-5}$ of the reference $0.23122$ value.  The remaining
non-canonical step is to derive the twisted-threshold beta coefficients and
$\alpha_{\rm UBT}^{-1}\simeq24.3$ directly from the UBT compact odd-spinor
spectrum.
\end{abstract}

\begin{statusbox}
\textbf{Boundary theorem:} $\sin^2\theta_W(M_{\rm UBT})=3/8$ [L1].\\
\textbf{Minimal SM one-loop descent:} obstruction; gives the wrong low-energy
value from the boundary.\\
\textbf{Twisted odd-spinor branch:} viable; one-loop RG gives
$\sin^2\theta_W(M_Z)=0.23114\ldots$ for $M_{\rm UBT}=2\times10^{16}$ GeV and
$\alpha_{\rm UBT}^{-1}=24.3$.\\
\textbf{Remaining gap:} derive the twisted Layer2 threshold spectrum and
unified coupling from $S[\Theta]$ rather than importing MSSM-like beta
coefficients.
\end{statusbox}

\section{Boundary condition}

The UBT hypercharge normalisation theorem gives
\[
  \sum_i Y_i^2\dim(r_i)=\frac{10}{3},
  \qquad
  \sum_i T_{3,i}^2\dim(r_i)=2.
\]
Thus
\[
  k_Y=\frac{5}{3},
  \qquad
  \frac{g_Y^2}{g_2^2}=\frac{3}{5},
\]
and therefore
\[
  \sin^2\theta_W(M_{\rm UBT})
  =
  \frac{g_Y^2}{g_2^2+g_Y^2}
  =
  \frac{3/5}{1+3/5}
  =
  \frac38.
\]

\section{One-loop RG equations}

For GUT-normalised couplings $\alpha_1$ and $\alpha_2$ we use
\[
  \frac{d\alpha_i^{-1}}{d\ln\mu}
  =-\frac{b_i}{2\pi},
\]
so that running from $M_{\rm UBT}$ down to $M_Z$ gives
\[
  \alpha_i^{-1}(M_Z)
  =
  \alpha_{\rm UBT}^{-1}
  +
  \frac{b_i}{2\pi}
  \ln\frac{M_{\rm UBT}}{M_Z}.
\]
The physical hypercharge coupling is
\[
  \alpha_Y=\frac35\alpha_1,
\]
and therefore
\[
  \sin^2\theta_W(M_Z)
  =
  \frac{\alpha_Y(M_Z)}{\alpha_Y(M_Z)+\alpha_2(M_Z)}.
\]

\section{Minimal SM branch: obstruction}

With minimal non-supersymmetric SM one-loop beta coefficients,
\[
  b_1=\frac{41}{10},
  \qquad
  b_2=-\frac{19}{6},
\]
and with representative values
\[
  M_{\rm UBT}=2\times10^{16}\,{\rm GeV},
  \qquad
  \alpha_{\rm UBT}^{-1}=40,
\]
one obtains
\[
  \sin^2\theta_W(M_Z)=0.18547\ldots.
\]
Thus the minimal SM one-loop branch is not sufficient to connect the UBT
$3/8$ boundary to the observed $Z$-pole value.

\begin{statusbox}
This is an obstruction result, not a failure of the boundary theorem.  It says
that the low-energy Weinberg angle requires threshold structure beyond minimal
non-supersymmetric one-loop running if the boundary is placed near
$2\times10^{16}$ GeV.
\end{statusbox}

\section{Twisted odd-spinor branch}

The alpha proof uses the canonical bridge that odd winding spinor modes carry
Berezin/Grassmann measure.  The same sector naturally suggests a twisted
Layer2 threshold in which the active electroweak compact spectrum has
MSSM-like one-loop coefficients
\[
  b_1=\frac{33}{5},
  \qquad
  b_2=1.
\]
This should not be read as importing a full MSSM model.  It is a compact way of
recording the effective one-loop beta coefficients expected from an odd-spinor
paired threshold spectrum.  The theorem still required is to derive this
spectrum directly from the UBT Layer2 Hilbert space.

Using
\[
  M_{\rm UBT}=2\times10^{16}\,{\rm GeV},
  \qquad
  M_Z=91.1876\,{\rm GeV},
\]
we have
\[
  L:=\frac{1}{2\pi}\ln\frac{M_{\rm UBT}}{M_Z}=5.2555492367\ldots.
\]
For $\alpha_{\rm UBT}^{-1}=24.3$,
\begin{align}
  \alpha_1^{-1}(M_Z)&=24.3+\frac{33}{5}L=58.98662496\ldots,\\
  \alpha_2^{-1}(M_Z)&=24.3+L=29.55554924\ldots.
\end{align}
Hence
\[
  \alpha_Y^{-1}(M_Z)=\frac53\alpha_1^{-1}(M_Z)=98.31104160\ldots,
\]
and
\[
  \boxed{
  \sin^2\theta_W(M_Z)=0.2311436399\ldots
  }.
\]
This differs from the reference value $0.23122$ by about
$-7.6\times10^{-5}$.

Solving the same one-loop twisted branch for the reference central value gives
\[
  \alpha_{\rm UBT}^{-1}=24.3254657224\ldots.
\]
At that point
\[
  \alpha_{\rm EM}^{-1}(M_Z)=127.934499434\ldots,
\]
which is in the expected $Z$-pole range.

\section{What is proven and what remains open}

The following chain is now explicit:
\[
  \text{UBT hypercharge norms}
  \Longrightarrow
  \sin^2\theta_W(M_{\rm UBT})=\frac38,
\]
then, on the twisted odd-spinor RG branch,
\[
  \frac38
  \xrightarrow[\alpha_{\rm UBT}^{-1}\simeq24.3]
  {b_1=33/5,\ b_2=1,\\ M_{\rm UBT}\simeq2\times10^{16}{\rm GeV}}
  \sin^2\theta_W(M_Z)=0.23114\ldots.
\]
The remaining proof tasks are precise:
\begin{enumerate}
  \item derive the twisted odd-spinor threshold spectrum and hence
  $b_1=33/5$, $b_2=1$ from the UBT Layer2 Hilbert space;
  \item derive $M_{\rm UBT}\simeq2\times10^{16}$ GeV from the compact
  $\psi$-radius or moduli-stabilisation sector;
  \item derive $\alpha_{\rm UBT}^{-1}\simeq24.3$ from the same normalisation
  rather than taking it as a representative unification value.
\end{enumerate}

\begin{statusbox}
\textbf{Current status.}  Weinberg angle: boundary theorem [L1]; RG-to-$M_Z$
viable on the twisted odd-spinor branch [L1 conditional].  The observed
low-energy value is not yet fully proven from canonical UBT because the twisted
threshold spectrum and scale still require derivation.
\end{statusbox}

\end{document}
